% Figure 1 (teaser). One agent session, recorded two ways.
\newcommand{\FigOneCaption}{One hour with a coding agent, recorded two ways. The
developer asked for three things and the agent did a fourth nobody asked for, and
the two records differ in whether anything in the repository says which is which.}

\newcommand{\FigOneDescription}{A three panel diagram. The left panel lists three
requests a developer typed to a coding agent, each marked with a small figure of a
person, and below them one change the agent made without being asked, marked with
a small figure of a robot. Each request has its own colour. The middle panel shows
two git commits drawn as cards containing file names with coloured diff hunks
under each, and hunks of every colour are mixed together inside both commits. The
right panel shows four horizontal rails, one per request, with square, circular and
triangular marks placed along each rail under the commit in which that edit
happened. Three marks on three different rails are ringed to show that all three
edited the same function. A strip along the bottom asks the same task of both
records: taking the caching back out means finding six of the fifteen hunks by hand
under git, and one named command under sgt.}

\begin{tikzpicture}[x=1mm, y=1mm, line cap=round, font=\sffamily]

% ===================== panel A: what the developer said ====================
\panel{0}{15}{36}{42}{sgtInk}{WHAT WAS ASKED}
\node[note, anchor=north west, text=sgtGrey] at (2.2,53.6) {one sitting, 14:02 to 15:20};

% #1 top y, #2 height, #3 colour, #4 wash, #5 time, #6 text
\newcommand{\reqcard}[6]{%
  \fill[#4, rounded corners=0.7mm] (2.2,{#1-#2}) rectangle (33.8,#1);
  \fill[#3, rounded corners=0.7mm] (2.2,{#1-#2}) rectangle (3.3,#1);
  \node[anchor=north west, inner sep=0, text=#3,
        font=\sffamily\bfseries\fontsize{4.9}{5.6}\selectfont] at (4.4,{#1-0.5}) {#5};
  \node[anchor=north west, text width=25.5mm, align=left, inner sep=0,
        font=\sffamily\fontsize{5.7}{6.5}\selectfont, text=sgtInk]
    at (4.4,{#1-2.7}) {#6};}

\reqcard{51.0}{7.0}{fRate}{wRate}{14:02}{``rate limit the public API, 60 a minute per key''}
\whodev{31.6}{47.5}{fRate}
\reqcard{42.7}{7.0}{fCache}{wCache}{14:31}{``the price lookup is slow, cache it''}
\whodev{31.6}{39.2}{fCache}
\reqcard{34.4}{7.0}{fRetry}{wRetry}{15:01}{``retry the upstream call on a 5xx''}
\whodev{31.6}{30.9}{fRetry}
\reqcard{26.1}{10.0}{sgtGrey}{wGrey}{NOT ASKED FOR}{the agent renamed \fsym{fetch} to \fsym{fetch\_row} and fixed its four callers}
\whoagent{31.6}{21.6}{sgtGrey}

% ===================== panel B: what git stored ============================
\panel{39}{15}{63}{42}{sgtGrey}{WHAT GIT STORED}
\node[note, anchor=north west, text=sgtGrey] at (41.2,53.6) {2 commits, 5 files, 15 hunks};

% #1 x, #2 sha, #3 message
\newcommand{\commitblock}[3]{%
  \draw[rounded corners=0.8mm, line width=0.3pt, color=sgtInk!18, fill=white]
    (#1,18.7) rectangle ({#1+28.4},51.2);
  \node[anchor=north west, inner sep=0, text=sgtInk,
        font=\ttfamily\fontsize{5.2}{6}\selectfont] at ({#1+1.4},50.5) {#2};
  \node[anchor=north west, text width=24mm, align=left, inner sep=0, text=sgtInk!70,
        font=\sffamily\fontsize{5.3}{6.1}\selectfont] at ({#1+1.4},47.7) {#3};}
% #1 x, #2 y, #3 name
\newcommand{\fname}[3]{\node[anchor=north west, inner sep=0, text=sgtInk!80,
  font=\ttfamily\fontsize{5.2}{6}\selectfont] at ({#1+1.4},#2) {#3};}

\commitblock{41.2}{c39ab2}{``add rate limiting and price cache''}
\fname{41.2}{42.0}{api.py}
  \hunk{43.4}{38.8}{15}{fRate}\hunk{43.4}{37.5}{10}{fRate}
  \hunk{43.4}{36.2}{13}{fRate}\hunk{43.4}{34.9}{18}{fCache}
\fname{41.2}{33.3}{cache.py}
  \hunk{43.4}{30.1}{21}{fCache}\hunk{43.4}{28.8}{13}{fCache}
\fname{41.2}{27.2}{db.py}
  \hunk{43.4}{24.0}{11}{fCache}
\fname{41.2}{22.4}{util.py}
  \hunk{43.4}{19.2}{9}{sgtGrey}

\commitblock{72.4}{7f10de}{``retries, plus fixes to the cache''}
\fname{72.4}{42.0}{api.py}
  \hunk{74.6}{38.8}{12}{fRetry}\hunk{74.6}{37.5}{17}{fCache}
\fname{72.4}{35.9}{retry.py}
  \hunk{74.6}{32.7}{20}{fRetry}\hunk{74.6}{31.4}{14}{fRetry}\hunk{74.6}{30.1}{18}{fRetry}
\fname{72.4}{28.5}{cache.py}
  \hunk{74.6}{25.3}{10}{fCache}
\fname{72.4}{23.7}{util.py}
  \hunk{74.6}{20.5}{15}{sgtGrey}

% the colours in this panel are the paper's, not git's, and saying so is the point
\node[anchor=north west, inner sep=0, text=sgtInk!70,
      font=\sffamily\fontsize{5.2}{6}\selectfont] at (41.2,18.1)
  {We coloured these hunks. Git stores no such link.};

% ===================== panel C: what sgt stored ===========================
\panel{105}{15}{69}{42}{sgtInk}{WHAT \textsf{sgt} STORED}

\draw[commitcol] (128.5,27.2) -- (128.5,51.6);
\draw[commitcol] (150.0,27.2) -- (150.0,51.6);
\node[note, text=sgtGrey, anchor=south] at (128.5,51.8) {c39ab2};
\node[note, text=sgtGrey, anchor=south] at (150.0,51.8) {7f10de};

\newcommand{\raillabel}[3]{\node[lbl, anchor=east, text=#3] at (123.6,#1) {#2};}
% A ring saying: this mark and the other ringed marks are the same function.
% Solid, and drawn solid on purpose. The legend two rows below this one gives a
% dotted outline to an edit that is no longer included, and every other figure
% spends dotted the same way, so a dotted ring here would mark the one fact in
% the panel that \sgt{} does record as the one thing it does not.
\newcommand{\shared}[2]{\draw[color=sgtInk!50, line width=0.3pt]
  (#1,#2) circle (2.0);}

% rate limiting
\draw[rail=fRate] (125.4,48.5) -- (171.6,48.5);
  \raillabel{48.5}{rate limiting}{fRate}
\opadd{128.9}{48.5}{fRate}\opadd{131.8}{48.5}{fRate}\opadd{134.7}{48.5}{fRate}
\opext{137.6}{48.5}{fRate}\ckpt{140.0}{48.5}{fRate}
\shared{137.6}{48.5}

% price cache
\draw[rail=fCache] (125.4,42.0) -- (171.6,42.0);
  \raillabel{42.0}{price cache}{fCache}
\opadd{128.9}{42.0}{fCache}\opadd{131.8}{42.0}{fCache}\opext{134.7}{42.0}{fCache}
\opext{137.6}{42.0}{fCache}\ckpt{140.0}{42.0}{fCache}
\oprew{150.4}{42.0}{fCache}\opext{153.3}{42.0}{fCache}\ckpt{155.7}{42.0}{fCache}
\shared{137.6}{42.0}

% retry policy
\draw[rail=fRetry] (125.4,35.5) -- (171.6,35.5);
  \raillabel{35.5}{retry policy}{fRetry}
\opadd{150.4}{35.5}{fRetry}\opadd{153.3}{35.5}{fRetry}\opext{156.2}{35.5}{fRetry}
\ckpt{158.6}{35.5}{fRetry}
\shared{156.2}{35.5}

% the rename nobody asked for
\draw[rail=sgtGrey] (125.4,29.0) -- (171.6,29.0);
  \raillabel{29.0}{row accessors}{sgtGrey}
\opmove{128.9}{29.0}{sgtGrey}\opext{131.8}{29.0}{sgtGrey}\opext{134.7}{29.0}{sgtGrey}
\opext{137.6}{29.0}{sgtGrey}\opext{150.4}{29.0}{sgtGrey}
\whoagent{106.4}{29.0}{sgtGrey}

% What the rings mean. The leader stops on the ring's own circumference rather
% than short of it, because the checkpoint tick stands where it used to land and
% an arrow between a ring and a tick names whichever of the two the reader
% happens to look at.
\annot{142.8}{45.4}{139.3}{43.1}{20}{sgtInk}
\node[anchor=west, text width=25mm, align=left, inner sep=0, text=sgtInk!75,
      font=\sffamily\fontsize{5.1}{5.9}\selectfont] at (143.2,45.4)
  {\fsym{api.py::handle}, written by three requests};

\node[note, anchor=east, text=sgtGrey] at (171.8,25.5) {18 edits, 14 functions, 5 files};

% legend, inside the panel whose marks it explains
\newcommand{\leg}[3]{\node[anchor=west, inner sep=0, text=sgtInk!70,
  font=\sffamily\fontsize{5.1}{5.9}\selectfont] at ({#1+2.4},#2) {#3};}
\opadd{107.6}{21.6}{sgtInk!55}\leg{107.6}{21.6}{new}
\opext{119.5}{21.6}{sgtInk!55}\leg{119.5}{21.6}{extended}
\oprew{137.5}{21.6}{sgtInk!55}\leg{137.5}{21.6}{reworked}
\opmove{156.0}{21.6}{sgtInk!55}\leg{156.0}{21.6}{renamed}
\ckpt{107.6}{18.2}{sgtInk!55}\leg{107.4}{18.2}{checkpoint}
\opghost{127.5}{18.2}{sgtInk!55}\leg{127.5}{18.2}{taken out}
\whodev{145.5}{18.2}{sgtInk!55}\leg{145.1}{18.2}{developer}
\whoagent{161.5}{18.2}{sgtInk!55}\leg{161.1}{18.2}{agent}

% ===================== the same task, of both records =====================
\plainpanel{0}{0}{174}{12.6}{wGrey}

\node[anchor=north west, text width=32mm, align=left, inner sep=0, text=sgtInk,
      font=\sffamily\bfseries\fontsize{6.0}{7.0}\selectfont] at (2.2,10.6)
  {A week later, the price cache has to come back out.};

\node[anchor=north west, text width=59mm, align=left, inner sep=0, text=sgtInk!85,
      font=\sffamily\fontsize{5.7}{6.7}\selectfont] at (41.2,10.6)
  {Find 6 of the 15 hunks, in 3 files, in 2 commits. Split both commits by hand,
   then check the \fsym{fetch\_row} rename did not ride along.};

\node[anchor=north west, inner sep=0, text=sgtInk,
      font=\ttfamily\fontsize{5.9}{7}\selectfont] at (107.2,10.6)
  {\$ sgt revert "price cache"};
\node[anchor=north west, inner sep=0, text=sgtInk!75,
      font=\sffamily\fontsize{5.5}{6.5}\selectfont] at (107.2,6.6)
  {6 edits, 4 functions, 3 files, one name};

\end{tikzpicture}
